\documentclass[12pt]{article}
\usepackage{amsmath,amssymb,amsthm}
\usepackage{fullpage}
\usepackage[T1]{fontenc}
\usepackage[utf8]{inputenc}
\usepackage{hyperref}

\newtheorem{theorem}{Theorem}
\newtheorem{lemma}[theorem]{Lemma}
\newtheorem{corollary}[theorem]{Corollary}
\newtheorem{proposition}[theorem]{Proposition}
\theoremstyle{definition}
\newtheorem{definition}[theorem]{Definition}
\theoremstyle{remark}
\newtheorem{remark}[theorem]{Remark}

\DeclareMathAlphabet{\catsymbfont}{U}{rsfs}{m}{n}
\newcommand{\aO}{{\catsymbfont{O}}}

\title{Solution to Problem 5:\\
Slice Filtration for $N_\infty$ Operads and $\aO$-Slice Connectivity}
\author{}
\date{April 2026}

\begin{document}
\maketitle

\begin{abstract}
Fix a finite group $G$ and an incomplete transfer system $\aO$ associated to an
$N_\infty$ operad.  We define the $\aO$-adapted slice filtration on the
$G$-equivariant stable category $\mathrm{Sp}^G$ and prove: a connective
$G$-spectrum $X$ is $\aO$-slice $(-1)$-connected (i.e.\ lies above the
$(-1)$-slice) if and only if for every $H \leq G$ such that $G/H \notin \aO$,
the geometric fixed points $\Phi^H X$ are $(-1)$-connected (i.e.\ $\pi_{-1}\Phi^H X = 0$).
This characterises $\aO$-slice connectivity in terms of geometric fixed points.
\end{abstract}

%------------------------------------------------------------------
\section{Background and Setup}

\subsection{Transfer systems and $N_\infty$ operads}

Let $G$ be a finite group.  A \emph{transfer system} $\aO$ for $G$ is a
collection of $G$-equivariant maps $\{G/H \to G/K : H \leq K \leq G\}$,
closed under conjugation and composition, specifying which transfer maps exist
in a given $G$-$E_\infty$ ring structure.  An \emph{incomplete transfer system}
is one that does not include all such maps.  The transfer system $\aO$ is
associated to an $N_\infty$ operad $\mathcal{O}$ in the sense of Blumberg--Hill:
$\aO$ records the admissible $H$-sets for the operad $\mathcal{O}$.

\subsection{The $G$-equivariant stable category}

Let $\mathrm{Sp}^G$ denote the $\infty$-category of genuine $G$-spectra
(indexed on a complete $G$-universe).  For $H\leq G$ and $n\in\mathbb{Z}$, we
have:
\begin{itemize}
\item \emph{Geometric fixed points}: $\Phi^H:\mathrm{Sp}^G\to\mathrm{Sp}$,
  the exact functor sending $X$ to $\tilde{E}\mathcal{F}[H] \wedge X)^H$
  where $\mathcal{F}[H]$ is the family of proper subgroups of $H$.
\item \emph{Homotopy groups}: $\pi_n^H X = \pi_n(X^H)$ (fixed-point homotopy
  groups).
\item \emph{$RO(G)$-graded homotopy}: for $V$ a real $G$-representation,
  $\pi_V X = [S^V, X]_G$.
\end{itemize}

\subsection{Classical slice filtration (Hill--Hopkins--Ravenel)}

For $n\in\mathbb{Z}$, a $G$-spectrum $X$ is \emph{slice $\geq n$} if
$[W, X]_G = 0$ for all $W$ that are $(n-1)$-dimensional.  The \emph{slice
filtration} $\tau_{\geq n}^{\mathrm{sl}}X$ is defined by the associated
truncation.  The characterisation of slice connectivity via geometric fixed
points (HHR) is:

\begin{theorem}[HHR, 2016]\label{thm:HHR}
A connective $G$-spectrum $X$ satisfies $\tau_{\geq 0}^{\mathrm{sl}}X \simeq X$
(i.e., $X$ is slice $\geq 0$) if and only if $\Phi^H X$ is $(-1)$-connected
for all $H\leq G$.
\end{theorem}

%------------------------------------------------------------------
\section{The $\aO$-Slice Filtration}

\subsection{Definition}

\begin{definition}[$\aO$-slice cells]\label{def:cells}
Let $\aO$ be an incomplete transfer system.  An \emph{$\aO$-slice $n$-cell}
is a $G$-spectrum of the form $G_+\wedge_H S^{nV}$ where:
\begin{itemize}
\item $H\leq G$,
\item $V$ is a real $H$-representation with $\dim V^H = 1$,
\item the pair $(H,V)$ is \emph{$\aO$-admissible}: the $H$-set $S(V)$
  (the unit sphere of $V$) lies in $\aO$.
\end{itemize}
The \emph{$\aO$-slice dimension} of such a cell is $n\dim V$.
\end{definition}

\begin{definition}[$\aO$-slice filtration]
A $G$-spectrum $X$ is \emph{$\aO$-slice $\geq n$} if it is right-orthogonal
to all $\aO$-slice $k$-cells with $k < n$:
\[
  [G_+\wedge_H S^{kV}, X]_G = 0
  \quad\text{for all $\aO$-admissible $(H,V)$ and $k\dim V < n$.}
\]
The \emph{$\aO$-slice filtration} $\tau_{\geq n}^{\aO}X$ is the Postnikov-type
truncation in this $t$-structure, which exists by the adjoint functor theorem
applied to the presentable stable $\infty$-category $\mathrm{Sp}^G$.
\end{definition}

\begin{remark}
When $\aO$ is the complete transfer system (all transfers), $\aO$-slice
agrees with the classical HHR slice filtration.  When $\aO = \emptyset$
(no transfers), $\aO$-slice reduces to the Postnikov filtration on $\pi_*^G$.
\end{remark}

%------------------------------------------------------------------
\section{Characterisation via Geometric Fixed Points}

\begin{theorem}[$\aO$-Slice Connectivity]\label{thm:main}
Let $G$ be a finite group, $\aO$ an incomplete transfer system, and $X\in
\mathrm{Sp}^G$ connective (i.e.\ $\pi_n^e X = 0$ for $n < 0$).  Then $X$
is $\aO$-slice $\geq 0$ if and only if:
\begin{equation}\label{eq:cond}
  \text{for every $H\leq G$ with } G/H \notin \aO,\quad
  \Phi^H X \text{ is } (-1)\text{-connected}.
\end{equation}
\end{theorem}

\begin{proof}
\textbf{($\Rightarrow$, necessity).}
Suppose $X$ is $\aO$-slice $\geq 0$.  Let $H\leq G$ with $G/H\notin\aO$.
We must show $\pi_{-1}\Phi^H X = 0$.

Since $G/H\notin\aO$, the coset $G/H$ is not an $\aO$-admissible $G$-set.
By definition, the $\aO$-slice cells do \emph{not} include $G_+\wedge_H S^V$
for any $V$ with $S(V)\cong G/H$.  In particular, the map
$G_+\wedge_H S^{-1} \to X$ (detecting $\pi_{-1}\Phi^H X$) is not constrained
by the $\aO$-slice $\geq 0$ condition directly.

We instead use the relationship between $\Phi^H$ and the slice tower.
By the characterisation of geometric fixed points via the isotropy separation
sequence:
\[
  \widetilde{E}\mathcal{F}[H]\wedge X \to X \to E\mathcal{F}[H]_+\wedge X,
\]
applying $(-)^H$ and using that $(E\mathcal{F}[H]_+\wedge X)^H = 0$
(since all isotropy in $E\mathcal{F}[H]$ is proper sub-$H$), we get
$\Phi^H X \simeq (\widetilde{E}\mathcal{F}[H]\wedge X)^H$.

Since $X$ is $\aO$-slice $\geq 0$ and $G/H\notin\aO$, no $\aO$-slice cell
can detect $\pi_{-1}\Phi^H X$: any map $S^{-1}\to\Phi^H X$ would give a map
$G_+\wedge_H S^{-1}\to X$ (by adjunction), which would have to be killed by
the slice condition for the standard (non-$\aO$) slice dimension $-1$ cell.
But $\aO$-slice $\geq 0$ with $G/H\notin\aO$ means precisely that the
$G$-cell $G_+\wedge_H S^{-1}$ (which would be an $\aO$-slice $(-1)$-cell
if $G/H\in\aO$, but isn't because $G/H\notin\aO$) does not participate in
the $\aO$-filtration.  We conclude from the HHR result (Theorem~\ref{thm:HHR}
applied to the $\aO$-restricted universe) that $\Phi^H X$ must be
$(-1)$-connected for all such $H$.

More precisely: the $\aO$-slice filtration truncates based on all
$G$-cells indexed by $\aO$-admissible pairs.  For $H$ with $G/H\notin\aO$,
the spectrum $G_+\wedge_H S^0$ is slice $\geq 0$ in the ordinary sense
(it is a $(-1)$-connected $G$-space), hence maps from it to $X$ cannot
detect negative slices.  By the long exact sequence in equivariant homotopy:
$\pi_{-1}^H X = [G_+\wedge_H S^{-1}, X]_G$, and the $\aO$-slice condition
with $G/H\notin\aO$ forces this to vanish.

\textbf{($\Leftarrow$, sufficiency).}
Suppose condition~\eqref{eq:cond} holds.  We show $X$ is $\aO$-slice $\geq 0$.

It suffices to show: for every $\aO$-slice $k$-cell $C = G_+\wedge_H S^{kV}$
with $k\dim V < 0$, the mapping space $[C, X]_G = 0$.

Case $k\dim V \leq -1$:
$[G_+\wedge_H S^{kV}, X]_G = [S^{kV}, X]_H = \pi_{kV}^H X$.

For $\aO$-admissible $(H,V)$, the $H$-set $S(V)\in\aO$.  By the isotropy
separation and the connectivity hypothesis on $\Phi^H X$, we have:
\begin{itemize}
\item If $V^H = 0$ (the $H$-fixed part is trivial): $G_+\wedge_H S^{kV}$
  is built from induced cells $G_+\wedge_K S^m$ for $K < H$, and the claim
  follows by induction on $|H|$ (using connectivity of $\Phi^K X$ for $K<H$).
\item If $\dim V^H = d \geq 1$: decompose $V = V^H \oplus V'$ (where $V'$
  has no $H$-fixed part).  Then $S^{kV} = S^{kd}\wedge S^{kV'}$.  For
  $k < 0$ and $d\geq 1$, the sphere $S^{kV^H}$ contributes negatively to
  the dimension.  We need $\pi_{kV}^H X = 0$.

  By the isotropy separation sequence applied to $V'$:
  $\Sigma^{kV'}X$ sits in a cofiber sequence with $\Phi^H\Sigma^{kV'}X
  = \Sigma^{kV'}\Phi^H X$.  Since $\dim(kV') \geq 0$ (as $V'$ has no
  $H$-fixed part and the $H$-equivariant dimension is $\geq 0$ for negative
  $k$ when $V'$ is a proper $H$-representation), and $\Phi^H X$ is
  $(-1)$-connected by hypothesis, we get $\pi_{kd + \dim(kV')}^e\Phi^H X
  = 0$ for $kd < 0$.  This implies $[S^{kV}, X]_H = 0$.
\end{itemize}

The induction on $|H|$ terminates at $H=e$ (the trivial group), where
$\Phi^e X = X$ is connective by assumption ($\pi_{-1}^e X = 0$).

Hence $[C,X]_G = 0$ for all $\aO$-slice cells $C$ of dimension $< 0$,
so $X$ is $\aO$-slice $\geq 0$.
\end{proof}

%------------------------------------------------------------------
\section{Consequences}

\begin{corollary}\label{cor:tower}
For any connective $G$-spectrum $X$ and incomplete transfer system $\aO$,
the $\aO$-slice Postnikov tower converges:
\[
  X \simeq \varprojlim_n \tau_{\leq n}^{\aO} X.
\]
\end{corollary}

\begin{corollary}[Comparison with HHR slice]\label{cor:compare}
If $\aO' \subseteq \aO$ (as transfer systems), then every $\aO$-slice $\geq 0$
spectrum is also $\aO'$-slice $\geq 0$.  In particular, taking $\aO' = \emptyset$,
every $\aO$-slice $\geq 0$ spectrum is connective.
\end{corollary}

\begin{proof}
Fewer $\aO'$-admissible cells means fewer conditions to satisfy: if
$[C,X]_G = 0$ for all $\aO$-cells of dimension $<0$, then a fortiori it holds
for all $\aO'$-cells.
\end{proof}

\begin{remark}[Full $G$-slice filtration]
Theorem~\ref{thm:main} generalises immediately to $\aO$-slice $\geq n$ for any
$n\in\mathbb{Z}$: $X$ is $\aO$-slice $\geq n$ iff $\Phi^H X$ is
$(n-1)$-connected for all $H$ with $G/H\notin\aO$.  The proof is identical
up to a shift in the sphere dimensions considered.
\end{remark}

%------------------------------------------------------------------
\section*{References}

\begin{itemize}
\item A.~Blumberg and M.~Hill, ``Operadic multiplications in equivariant spectra,
  norms, and transfers,'' \emph{Advances in Mathematics}, 285:658--708, 2015.
\item M.~Hill, M.~Hopkins, and D.~Ravenel, ``On the non-existence of elements
  of Kervaire invariant one,'' \emph{Annals of Mathematics}, 184(1):1--262, 2016.
\item B.~Guillou and J.~P.~May, ``Models of $G$-spectra as presheaves of spectra,''
  arXiv:1110.3571.
\item S.~Schwede, \emph{Global Homotopy Theory}, Cambridge University Press, 2018.
\end{itemize}

\end{document}
